\documentclass[10pt]{article}
\usepackage{pennysia}

\hypersetup{
  pdftitle={Pennysia: An Open Liquidity Network},
  pdfauthor={Pennysia, Inc.},
  pdfsubject={Whitepaper v1.0},
  pdfkeywords={Pennysia, liquidity, DeFi, chain abstraction, intent, exchange}
}

\title{\textbf{Pennysia: An Open Liquidity Network}}
\author{Pennysia, Inc.}
\date{October 2026}

\begin{document}
\maketitle

\begin{abstract}
\noindent% !TEX root = ../main.tex
Pennysia is an open liquidity network, a single exchange layer where every source of liquidity, wherever it sits, competes for the user's order. Decentralized exchanges, aggregators, intent solvers, market makers, perpetual and prediction protocols, vaults, fiat ramps, and licensed venues all answer the same auction. The best executable price wins. The user signs once. The network routes, executes, and verifies the fill. For builders, a venue connects once and competes for every order from day one.

\noindent Pennysia aims to combine the performance, liquidity access, and usability of centralized exchanges (CEXs) with the self-custody of decentralized exchanges (DEXs). The complexity of wallet management, gas payments, and chain routing is abstracted away, while users retain their choice of wallet and chain. Fiat enters and exits on card and bank rails. Liquidity is organized into two lanes: a \textbf{permissionless lane} of open venues, onchain and offchain, open to everyone, and a \textbf{regulated lane} of gated liquidity (centralized market makers, licensed venues, compliant pools, fiat ramps) unlocked by optional verification.

\noindent An experimental deployment on Ethereum proved the model: 30 venues competing in parallel per swap, atomic in-transaction fallback, and a settlement layer that enforces signed terms on every order. Pennysia v1, the Open Liquidity Network, generalizes that proof to every chain and every service class, a single interface for the world's liquidity.

\end{abstract}

\vspace{0.5em}
{\small\noindent\textit{This paper describes Pennysia v1, the Open Liquidity Network, as designed. Design decisions open at publication are noted where they occur; the Ethereum deployment documented in Appendix~\ref{app:v1} demonstrates the core mechanisms in production.}}

% !TEX root = ../main.tex
\section{Introduction}\label{sec:introduction}

Crypto liquidity is fragmented along two axes at once~\cite{sok-defi}, and every exchange today serves only a slice of the result.

The first axis is \textbf{chains}. Liquidity is spread across dozens of networks (Ethereum and its L2s, Solana, BNB Chain, and a growing set of new execution environments), each with its own DEXs, aggregators, market makers, and intent protocols. No single venue wins on every pair, size, and moment. A trader who chases prices across chains inherits bridge risk~\cite{bridges}, gas-token juggling, and a different mental model per network. Chains behave like disconnected internets. A service lives where it deployed, and reaching it elsewhere means switching providers. Even the strongest onchain venues remain single-domain: excellent execution on their own chains and instruments, with the rest of the market outside reach.

The second axis is the \textbf{protocol boundary}: the split between \emph{open} and \emph{closed} protocols. Open protocols (DEXs, intent networks, onchain markets and vaults) are permissionless, transparent, and self-settling, but shallow on size~\cite{dexdepth}. Closed protocols (centralized books, market makers, OTC desks) hold deeper, tighter liquidity behind accounts, KYC, and custodial terms. Traders choose between two products and two trust models, with no common interface. The deepest quote in the world is invisible to an open-protocol trader.

Between the trader and both axes sits \textbf{the UX wall}. Self-custody is the founding premise of this space: Bitcoin was built so that holding and moving value requires no trusted intermediary~\cite{bitcoin}. Its cost has never been ideological; it is operational. Seed phrases gate the first ten minutes. Gas requires a native token on every chain. Chain selection, bridging, and per-venue approvals turn one trade into a checklist, and the endpoints (buying in with a card, cashing out to a bank account) still route through custodial platforms. Each step is a point where users exit to a centralized exchange's account login. The market's trade-off today: \textbf{easy to use, or yours.} Centralized exchanges offer the experience; decentralized venues offer the custody. No venue offers both. Figure~\ref{fig:journey} lays out both journeys.

\begin{figure}[ht]
\centering
\begin{tikzpicture}[font=\footnotesize, every node/.append style={text=chrome}]
\pgfmathsetlengthmacro{\colw}{0.94\textwidth/4}
\pgfmathsetlengthmacro{\labx}{0.12cm}
% frame: DeFi row 1.6--4.2, CeFi row 0--1.6
\draw[olline] (0,0) rectangle (0.94\textwidth,4.2);
\draw[olline] (0,1.6) -- (0.94\textwidth,1.6);
\draw[olline] (1*\colw,0) -- (1*\colw,4.2);
\draw[olline] (2*\colw,0) -- (2*\colw,4.2);
\draw[olline] (3*\colw,0) -- (3*\colw,4.2);
% column headers
\node[ollabel, anchor=south] at (0.5*\colw,4.2) {\footnotesize On-ramp};
\node[ollabel, anchor=south] at (1.5*\colw,4.2) {\footnotesize Wallet};
\node[ollabel, anchor=south] at (2.5*\colw,4.2) {\footnotesize Activity};
\node[ollabel, anchor=south] at (3.5*\colw,4.2) {\footnotesize Off-ramp};
% row labels
\node[ollabel, rotate=90] at (-0.4,2.9) {\footnotesize DeFi};
\node[ollabel, rotate=90] at (-0.4,0.8) {\footnotesize CeFi};
% DeFi row
\node[anchor=north west, align=left, text width=\colw-0.3cm] at (\labx,4.1)
  {1.~open a CEX or card-ramp\\ \hspace{0.6em}account\\ 2.~pass KYC either way\\ 3.~buy the asset\\ 4.~wait out the withdrawal lock};
\node[anchor=north west, align=left, text width=\colw-0.3cm] at (\colw+\labx,4.1)
  {5.~install a wallet\\ 6.~write down the seed phrase\\ 7.~acquire the gas token\\ 8.~withdraw to the right chain};
\node[anchor=north west, align=left, text width=\colw-0.3cm] at (2*\colw+\labx,4.1)
  {9.~pick venue and chain\\ 10.~approve token (tx + gas)\\ 11.~swap (gas)\\ 12.~bridge if needed (gas + wait)};
\node[anchor=north west, align=left, text width=\colw-0.3cm] at (3*\colw+\labx,4.1)
  {13.~sell into a stablecoin\\ 14.~send back to a CEX\\ 15.~sell for fiat\\ 16.~bank withdrawal};
% CeFi row
\node[anchor=north west, align=left, text width=\colw-0.3cm] at (\labx,1.5)
  {1.~open account\\ 2.~pass KYC\\ 3.~deposit by card or bank};
\node[anchor=north west, align=left, text width=\colw-0.3cm] at (\colw+\labx,1.5)
  {4.~balance and keys\\ \hspace{0.5em}held by the venue};
\node[anchor=north west, align=left, text width=\colw-0.3cm] at (2*\colw+\labx,1.5)
  {5.~trade on the venue's\\ \hspace{0.5em}internal ledger};
\node[anchor=north west, align=left, text width=\colw-0.3cm] at (3*\colw+\labx,1.5)
  {6.~bank withdrawal};
\end{tikzpicture}
\caption{Illustrative fiat-in to fiat-out journeys. Self-custody: ${\sim}16$ steps across three platforms, with seed custody, gas, and chain choice at every boundary; card funding still means a KYC'd ramp at a markup. Custodial: ${\sim}6$ steps, because the venue holds the keys. Design target: self-custody at the custodial step count.}
\label{fig:journey}
\end{figure}

Three enablers make this the moment to build it. Intent-based execution~\cite{intentauction} turned ``user signs, someone else transacts'' into production infrastructure, the mechanical basis for gasless, chainless trades. Embedded wallets matured~\cite{smartaccounts} to where a Web2 login (email, Google, Apple) provisions real EVM and Solana wallets, giving Web2 onboarding with user-side keys. And venue proliferation across both liquidity kinds makes comparison a routing problem, not a shopping problem. Figure~\ref{fig:dexshare} charts the migration already underway.

\begin{figure}[ht]
\centering
\begin{tikzpicture}
\begin{axis}[olchart,
  xmin=2021, xmax=2027,
  ymin=0, ymax=40,
  xtick={2021,2022,2023,2024,2025,2026},
  xticklabels={2021,2022,2023,2024,2025,2026},
  ylabel={DEX volume as \% of CEX spot volume},
  ytick={0,10,20,30,40},
]
% Monthly spot ratio, Jan 2021 -- Sep 2024 (CoinGecko).
\addplot[accent, thick] coordinates {
(2021.000,5.98) (2021.083,6.32) (2021.167,7.6) (2021.250,7.84) (2021.333,10.41) (2021.417,13.07) (2021.500,11.82) (2021.583,10.37) (2021.667,9.91) (2021.750,7.44) (2021.833,9.44) (2021.917,12.89) (2022.000,14.58) (2022.083,14.58) (2022.167,11.42) (2022.250,8.72) (2022.333,9.44) (2022.417,9.04) (2022.500,9.3) (2022.583,7.32) (2022.667,5.43) (2022.750,6.29) (2022.833,8.76) (2022.917,8.02) (2023.000,8.04) (2023.083,7.18) (2023.167,10.32) (2023.250,10.18) (2023.333,11.78) (2023.417,10.42) (2023.500,9.46) (2023.583,9.47) (2023.667,10.86) (2023.750,9.36) (2023.833,9.32) (2023.917,9.07) (2024.000,8.65) (2024.083,8.35) (2024.167,9.28) (2024.250,10.25) (2024.333,11.22) (2024.417,13.02) (2024.500,12.95) (2024.583,11.8) (2024.667,10.16)
};
% Reported landmark months: CoinGecko readings (circles),
% then The Block top-30-DEX series (square; wider venue set).
\addplot[accent, thick, mark=*, mark size=1.1pt] coordinates {
(2024.667,10.16) (2025.000,18.7) (2025.417,37.4) (2025.833,21.2)
};
\addplot[accent, thick, mark=square*, mark size=1.9pt] coordinates {
(2025.833,21.2) (2026.67,25.0)
};
\node[ollabel, anchor=south west, inner sep=2pt] at (axis cs:2026.67,25.0) {25\%};
\end{axis}
\end{tikzpicture}
\caption{DEX spot volume as a share of CEX spot volume, 2021--2026: 6\% to ${\sim}25\%$. Line and circles: CoinGecko~\cite{coingecko} (monthly through Sep.\ 2024; landmark months after). Square: The Block's wider top-30-DEX series, Sep.\ 2026~\cite{theblock}, indicative not strictly comparable. The migration is real but incomplete; the UX wall is the remaining ceiling.}
\label{fig:dexshare}
\end{figure}

Pennysia v1, the Open Liquidity Network (OLN), generalizes the Ethereum prototype's proof (Appendix~\ref{app:v1}) along both fragmentation axes: a unified layer spanning \textbf{every chain} and \textbf{every liquidity service} (spot today; perpetuals, prediction markets, and yield protocols next), wrapped in an account that removes the UX wall. Pennysia implements none of these services; it operates the network that connects to all of them. The end state is a superset of the centralized exchange: every essential feature (one login, fiat rails, deepest liquidity, more tokens, fast settlement) with custody that never leaves the user's keys. For builders the network is distribution: a new protocol connects once and competes for every order.


% !TEX root = ../main.tex
\section{Design Principles}\label{sec:principles}

Six commitments govern the design. Where an implementation decision remains open, the governing principle is stated in its place.

\begin{itemize}[leftmargin=1.5em]
  \item \textbf{Liquidity is venue-agnostic.} No venue is special. Every liquidity source and service protocol reduces to the same abstraction, \emph{a source that returns an executable price}. Pennysia builds no competing venue, pool, or protocol of its own. It operates the network that connects to all of them.

  \item \textbf{Non-custodial by construction.} The network never holds user funds. No pooled balances, no deposit accounts, no withdrawal queue. If the network disappeared tomorrow, every asset would remain in users' wallets.

  \item \textbf{Signature-first UX.} Users sign. The network transacts. A trade is one review and one signature. Gas, chain selection, submission, and confirmation are the network's problem.

  \item \textbf{Competition decides, not trust.} Quotes rank by \textbf{net executable output} (what lands in the wallet after gas, fees, and fill probability), not headline price. No venue buys placement. Listing is free, rank is earned per order.

  \item \textbf{Compliance is a lane, not a gate.} Verification is opt-in and rewarded with liquidity, never required as the price of entry.

  \item \textbf{Honest fees.} Costs live inside the quote, never after it. The network earns when execution beats the quote, not when it does not. There is no native token, and access is not gated by holdings.
\end{itemize}

% !TEX root = ../main.tex
\section{System Overview}\label{sec:architecture}

The architecture separates the user interface and account from quote selection and execution (Figure~\ref{fig:layers}).

\begin{figure}[ht]
\centering
\begin{tikzpicture}[font=\footnotesize, every node/.append style={text=chrome}]
  \draw[olline] (0,0) rectangle (0.97\textwidth,3.2);
  \draw[olline] (0,0.8) -- (0.97\textwidth,0.8);
  \draw[olline] (0,1.6) -- (0.97\textwidth,1.6);
  \draw[olline] (0,2.4) -- (0.97\textwidth,2.4);
  \draw[olline] (3.1,0) -- (3.1,3.2);
  \node[anchor=west] at (0.15,2.8) {\textbf{Experience}};
  \node[anchor=west, align=left] at (3.35,2.8) {one account, one balance, one ticket; chain \& gas hidden, status to finality};
  \node[anchor=west] at (0.15,2.0) {\textbf{Identity}};
  \node[anchor=west, align=left] at (3.35,2.0) {login $\rightarrow$ non-custodial wallets (embedded or external); optional KYC};
  \node[anchor=west] at (0.15,1.2) {\textbf{Quote market}};
  \node[anchor=west, align=left] at (3.35,1.2) {parallel fanout $\rightarrow$ ranked net output; permissionless lane $|$ regulated lane};
  \node[anchor=west] at (0.15,0.4) {\textbf{Execution}};
  \node[anchor=west, align=left] at (3.35,0.4) {signed order $\rightarrow$ best venue $\rightarrow$ verified fill; Settlement, intents, venue rails};
\end{tikzpicture}
\caption{Responsibilities of the four network layers.}
\label{fig:layers}
\end{figure}

Identity connects wallet access and verification to the account (Section~\ref{sec:account}). The quote market compares eligible providers (Section~\ref{sec:quote-market}), then the execution layer settles the signed order using the appropriate rail (Section~\ref{sec:execution}). The experience layer presents balances and order status throughout this process. Figure~\ref{fig:participants} shows how external participants connect to these layers.

\begin{figure}[ht]
\centering
\begin{tikzpicture}[font=\footnotesize, every node/.append style={text=chrome}]
  % Pennysia boundary
  \draw[olline, rounded corners=3pt] (3.2,-0.7) rectangle (11.6,2.5);
  \node[font=\footnotesize\bfseries] at (7.4,2.1) {Pennysia};
  % layers
  \node[olbox, minimum width=1.7cm] (exp) at (4.5,1.0) {Experience};
  \node[olbox, minimum width=1.7cm] (idn) at (6.7,1.0) {Identity};
  \node[olbox, minimum width=2.1cm] (qm)  at (9.0,1.0) {Quote market};
  \node[olbox, minimum width=1.8cm] (exe) at (10.4,-0.3) {Execution};
  % participants
  \node[olbox, ellipse, minimum width=2.0cm] (user) at (0.6,1.0) {\textbf{User}};
  \node[olbox, ellipse, minimum width=3.2cm, align=center] (wallet) at (5.5,-2.4)
    {\textbf{Wallet provider}\\ \scriptsize co-secures keys};
  \node[olbox, ellipse, minimum width=3.2cm, align=center] (verif) at (8.9,-2.4)
    {\textbf{KYC provider}\\ \scriptsize optional credential};
  \node[olbox, ellipse, minimum width=3.0cm, align=center, minimum height=2.4cm] (ven) at (14.5,0.6)
    {\textbf{Venues}\\ \scriptsize DEXs $\cdot$ intents\\ \scriptsize RFQ $\cdot$ regulated\\ \scriptsize ramps $\cdot$ vaults\\[2pt] \scriptsize\itshape ranked winner fills};
  % connections
  \draw[olarrow] (user.east) -- node[ollabel, above]{signs once} (exp.west);
  \draw[olarrow] (wallet.north) -- node[ollabel, left, pos=0.28]{key shares} (idn.south west);
  \draw[olarrow] (verif.north) -- node[ollabel, right, pos=0.28]{credential} (idn.south east);
  \draw[olarrow] ([yshift=5pt]qm.east) -- node[ollabel, above, pos=0.3]{orders} ([yshift=5pt]qm.east -| ven.west);
  \draw[olarrow] ([yshift=-5pt]qm.east -| ven.west) -- node[ollabel, below, pos=0.7]{quotes} ([yshift=-5pt]qm.east);
  \draw[olarrow] ([yshift=4pt]exe.east) -- node[ollabel, above]{execute} ([yshift=4pt]exe.east -| ven.west);
  \draw[olarrow] ([yshift=-4pt]exe.east -| ven.west) -- node[ollabel, below]{fill} ([yshift=-4pt]exe.east);
\end{tikzpicture}
\caption{Participant interfaces. Wallet and KYC providers connect through Identity; venues exchange quotes and settlement information with the quote market and execution layers. The selected venue receives the signed order.}
\label{fig:participants}
\end{figure}
\FloatBarrier

% !TEX root = ../main.tex
\section{The Account}\label{sec:account}

The account combines familiar login methods with user-controlled onchain wallets.

\subsection{Onboarding and custody}

A Web2 login (email, Google, or Apple) provisions \textbf{embedded wallets} for each supported chain family, initially EVM and Solana, without requiring a seed phrase or browser extension during signup. Assets reside at user-controlled onchain addresses, and moving them requires a wallet signature produced in the user's session. Table~\ref{tab:authority} defines each participant's authority.

\begin{table}[ht]
\centering
\caption{Authority boundaries: what each actor can and cannot do.}
\label{tab:authority}
\begin{tabular}{@{}l>{\raggedright\arraybackslash}p{0.37\textwidth}>{\raggedright\arraybackslash}p{0.35\textwidth}@{}}
\toprule
\textbf{Actor} & \textbf{Can do} & \textbf{Cannot do} \\
\midrule
User & Sign orders and transactions; export and self-manage keys & \emph{none} \\
Pennysia & Verify sessions, assemble quotes, submit signed orders, display state & Move user funds; sign for the user; recover wallets \\
Wallet provider & Co-secure key material (device/enclave shares); assist user-driven recovery & Unilaterally spend \\
\bottomrule
\end{tabular}
\end{table}

Session verification authorizes account operations and quote assembly. It does not grant Pennysia signing authority.

\subsection{Gas and wallet choice}

Executors pay gas and include its cost in the quote (Section~\ref{sec:execution}), so users can trade without acquiring a native gas token. Where a user must authorize an onchain action, such as an initial approval, the network sponsors or abstracts the gas.

External wallets, including extensions, WalletConnect connections, and hardware wallets, access the same quote market. The account provides a unified view and session for these wallets.

\subsection{In and out}

External ramp providers quote fiat-to-token and token-to-fiat conversions through the auction, with support for cards, bank transfers, Apple Pay, and Google Pay. These services sit in the regulated lane and require the applicable verification (Section~\ref{sec:credential}). Pennysia coordinates access without holding a custodial balance.

The wallet provider handles recovery through re-authentication and restoration of key shares. Users can export and manage their keys independently of Pennysia. The account layer supports externally owned accounts signing typed-data orders~\cite{eip712} and smart accounts~\cite{erc4337,eip7702} with session keys and programmable recovery; the order format in Appendix~\ref{app:order} is compatible with either.

% !TEX root = ../main.tex
\section{The Quote Market}\label{sec:quote-market}

Liquidity today sits on protocols the way accommodation sits on hotels. Thousands of providers serve the same essential need, at different prices and terms, and a user should no more browse every venue than a traveler should browse every hotel's website. The quote market is the routing layer. It turns a fragmented liquidity landscape into a single price. Every order is an auction, every eligible venue sees the same request simultaneously, and the best \emph{executable} answer wins~\cite{solverwelfare}.

\subsection{The auction}

A trade request (sell X, get Y) broadcasts to every venue in the user's enabled lanes. Formally, for user $u$ each venue $v$ in the eligible set $E(u)$ returns a quote $(y_v, c_v, p_v)$: gross output, all-in cost (gas, venue and protocol fees), and fill probability, with expected settlement latency folded into the ranking. Quotes answer within a response window $\tau$ set by the market. A venue that misses the window does not rank, so every quote in the auction is a live commitment rather than indicative pricing. The market selects the venue maximizing net executable output:

\begin{equation}
E(u) \;=\; \mathcal{V}_{\mathrm{perm}} \;\cup\; \mathbf{1}[\mathrm{verified}(u)]\cdot \mathcal{V}_{\mathrm{reg}},
\qquad
v^{*}(o) \;=\; \underset{v \in E(u)}{\arg\max}\; \underbrace{p_v \cdot (y_v - c_v)}_{\mathrm{net\; output}}
\end{equation}

\noindent A quote 0.3\% better that fails 5\% of the time loses to a slightly worse quote that settles. This ranking already runs on the Ethereum prototype; the OLN applies it identically across chains and venue classes.

\subsection{Two axes of liquidity}\label{sec:lanes}

Liquidity fragments along two independent axes: \textbf{permission} (who is allowed to serve a given user) and \textbf{location} (onchain protocols and contracts, or offchain books and inventory). The axes are orthogonal. A walletless provider quotes permissionlessly from offchain inventory, while a compliant pool is permissioned yet holds its liquidity onchain.

The permission axis defines the user's \emph{lane}. A permissionless venue is an open city, a regulated venue a licensed jurisdiction. The \textbf{permissionless lane} is open to every user: every open liquidity source (DEXs, aggregators, intent networks, RFQ makers, walletless providers), carrying the full prototype venue set. The \textbf{regulated lane} unlocks through optional verification (Section~\ref{sec:compliance}): centralized market makers, OTC desks, CEX-linked liquidity, and fiat ramps, where depth on large size tends to live.

The location axis predicts the execution rail (Section~\ref{sec:execution}). Atomic and intent rails settle onchain liquidity, and venue rails bridge offchain books to onchain delivery.

The lanes share one auction and one ranking. A regulated venue enters the same competition and wins by the same rule.

\subsection{Venue and service classes}\label{sec:venue-classes}

Figure~\ref{fig:axes} maps the landscape onto both axes.

\begin{figure}[ht]
\centering
\begin{tikzpicture}[font=\footnotesize, every node/.append style={text=chrome}]
  \pgfmathsetlengthmacro{\qw}{0.39\textwidth}
  \def\qh{2.3}
  % axes: x = permission, y = location
  \draw[olline] (-\qw,0) -- (\qw,0);
  \draw[olline] (0,-\qh) -- (0,\qh);
  % axis-end labels
  \node[font=\footnotesize\bfseries, anchor=east] at (-\qw,0) {Permissionless};
  \node[font=\footnotesize\bfseries, anchor=west] at (\qw,0) {Regulated};
  \node[font=\footnotesize\bfseries, anchor=south] at (0,\qh) {Onchain};
  \node[font=\footnotesize\bfseries, anchor=north] at (0,-\qh) {Offchain};
  % quadrants
  \node[align=center, text width=0.42\textwidth] at (-0.5*\qw,0.5*\qh)
    {\textbf{Open, self-settling}\\[0.3em]
     DEXs, aggregators, intents, RFQ makers\\[0.15em]
     perps, prediction markets, yield vaults};
  \node[align=center, text width=0.42\textwidth] at (0.5*\qw,0.5*\qh)
    {\textbf{Gated onchain}\\[0.3em]
     KYC'd pools, compliant vaults};
  \node[align=center, text width=0.42\textwidth] at (-0.5*\qw,-0.5*\qh)
    {\textbf{Open offchain}\\[0.3em]
     privacy-swap providers};
  \node[align=center, text width=0.42\textwidth] at (0.5*\qw,-0.5*\qh)
    {\textbf{Gated offchain}\\[0.3em]
     CEX makers, OTC desks, fiat ramps};
\end{tikzpicture}
\caption{Two axes of liquidity. The horizontal axis is permission, the user's lane; the vertical axis is where the liquidity lives. Every class enters the same auction behind the same signed-terms gate.}
\label{fig:axes}
\end{figure}

The classes extend beyond swaps. Any service returning an executable price for a position (a perpetual, a market share, a vault deposit, a fiat conversion) enters the same auction under the same abstraction. Coverage compounds with every connection. Each chain adds its venue set, each venue extends the token surface, and a token liquid anywhere becomes tradable everywhere.

Joining differs by lane. The permissionless lane is free and self-serve. Any venue that returns executable quotes against the network's adapter interface joins the auction, the model that grew the prototype to 30 venues. The regulated lane adds a compliance onboarding covering licenses, jurisdictions, and settlement rails. Once connected, a venue competes for every order from day one, ranked on output alone.

\subsection{Chainless quoting}

A quote request names assets and intent, not a chain. Chains are plumbing in this model, an internet provider for the services deployed on them, relevant to routing but irrelevant to the requester. The routing layer resolves where the fill happens (same-chain, cross-chain via fillers, or a regulated venue), and the user reviews net output and rail type, never a chain diagram. Chainless is the property of this liquidity network.

% !TEX root = ../main.tex
\section{Execution}\label{sec:execution}

Execution turns a signed order into settled tokens. The user's part ends at the signature.

\subsection{Signature-first order flow}\label{sec:signature-flow}

The user reviews a quote stating the worst acceptable output (\textbf{WAP}, the minimum received, net of all fees) and signs an order, a typed-data payload binding input, output, floor, deadline, and recipient (Appendix~\ref{app:order}). No transaction, no gas prompt, no chain selection. The network derives the full execution plan and carries it to completion.

\subsection{The three rails}\label{sec:rails}

Orders settle over three rails (Figure~\ref{fig:rails}).

\textbf{Rail 1: Atomic routes (same chain).} A Settlement-class router pulls the user's input once (signed allowance authorization), attempts ranked venues inside one transaction, first success wins. All routes fail $\rightarrow$ full revert, with no partial fills, no stranded funds, no second signature. This is the prototype's production model (Appendix~\ref{app:v1}), replicated per EVM chain.

\textbf{Rail 2: Intent rails (cross-chain / solver fills).} The signed order enters an intent settlement network. A filler commits to the quoted output, pays all gas, and delivers on the destination chain. User exposure is bounded by the signed terms: floor, deadline, recipient. Unfilled intents expire. An order never moves funds without delivery.

\textbf{Rail 3: Venue rails (regulated lane).} Centralized venues settle via escrow contract (atomic onchain exchange~\cite{atomicswaps}) or deposit-address rails, the privacy-swap pattern the prototype already runs, where input goes to a venue address and output arrives asynchronously. Both bounded by the signed terms and reconciled on completion. The counterparty model is disclosed per quote (Section~\ref{sec:risk}).

\begin{figure}[ht]
\centering
\begin{tikzpicture}[font=\footnotesize, every node/.append style={text=chrome}]
  \node[align=center] (ord) at (0,0) {\textbf{Signed order}\\ \scriptsize floor \,\textperiodcentered\, deadline \,\textperiodcentered\, recipient};
  \node[align=center] (fin) at (13.4,0) {\textbf{Verified fill}\\ \scriptsize or controlled failure};
  % three rail paths, each an arrow converging into the fill
  \draw[olarrow] (ord.east) -- (4.4,1.5) -- (10.2,1.5) -- (fin.150);
  \draw[olarrow] (ord.east) -- (fin.west);
  \draw[olarrow] (ord.east) -- (4.4,-1.5) -- (10.2,-1.5) -- (fin.210);
  % rail labels
  \node[anchor=south] at (7.3,1.55) {\textbf{Rail 1: Atomic} \textperiodcentered\ ranked venues, one transaction};
  \node[anchor=south] at (7.3,0.15) {\textbf{Rail 2: Intent} \textperiodcentered\ filler commits, pays gas};
  \node[anchor=south] at (7.3,-1.45) {\textbf{Rail 3: Venue} \textperiodcentered\ escrow, async};
\end{tikzpicture}
\caption{Three rails, one order format, one set of guarantees. The rails differ; the signed terms do not.}
\label{fig:rails}
\end{figure}

\subsection{Guarantees and failure}

One rule across rails: \textbf{funds move only under signed terms; a failed route returns control, not loss.} Formally, for every order $o$ signed by user $u$,

\begin{equation}
\mathrm{moved}(u, o) \;\Longrightarrow\; \mathrm{delivered}(u, o) \geq \mathrm{WAP}(o) \;\land\; t \leq \mathrm{deadline}(o)
\end{equation}

\noindent Atomic routes revert wholesale, intents expire unpaid, and venue-rail fills resolve to delivered or refunded. Where a route set exists, the network attempts the next ranked venue before surfacing failure. Table~\ref{tab:guarantees} summarizes the full contract.

\begin{table}[ht]
\centering\small
\caption{User-facing guarantees and the mechanism behind each.}
\label{tab:guarantees}
\begin{tabular}{@{}p{0.38\textwidth}p{0.52\textwidth}@{}}
\toprule
\textbf{Guarantee} & \textbf{Mechanism} \\
\midrule
Worst-case output known at signing & WAP embedded in the order; enforced per rail \\
No gas token required & Executor pays gas, priced into quote; sponsored where unavoidable \\
No chain management & Rail and destination are the network's problem \\
Bounded failure & Revert, expiry, or refund; never partial loss \\
Disclosed settlement time & Single-block same-chain; latency priced into the quote as ETA \\
One signature & Review $\rightarrow$ sign $\rightarrow$ status to finality \\
\bottomrule
\end{tabular}
\end{table}

Latency is part of the contract. Same-chain fills settle within a block, while slower rails carry a bounded ETA priced into the quote, targeting centralized-exchange responsiveness on the fast rail. The surface stays minimal. Pick input and output, see one net number and an ETA, sign once. The rail underneath is visible on demand, never required.

\begin{tbd}
Atomic routes as one Settlement-family contract per chain vs.\ a canonical intent-native executor shared across chains. Both preserve the semantics; the choice is engineering and audit scoping.
\end{tbd}

% !TEX root = ../main.tex
\section{Security Model}\label{sec:security}

The signed order is the enforcement boundary. Routes and quotes are validated against the signed payload before execution, and execution conditions are checked at submission because public-chain transaction ordering is not guaranteed~\cite{flashboys}. A fill is recorded only when delivery meets the signed minimum before the deadline to the specified recipient. Atomic routes spend through a single Settlement approval constrained by the authorized input, and controlled failure follows the procedure in Section~\ref{sec:rails}.

\subsection{Contract and venue exposure}

A unified approval concentrates contract risk in Settlement, making its implementation and verification the critical dependency. New contracts are intended to follow the prototype's process (Appendix~\ref{app:v1}): published source, invariant testing, and formal analysis where tractable, with deployment artifacts and verification results accompanying each chain rollout. Venue-side protections such as solver bonds or contract controls continue to apply underneath; Pennysia validates execution independently while retaining exposure to each venue's implementation and settlement model.

\subsection{Account security}

Table~\ref{tab:authority} separates session authority from signing authority. A compromised session exposes account operations but cannot move assets, because every transfer requires a wallet signature. Embedded wallets rely on co-secured key material and provider-assisted recovery~\cite{threshold}; users who prefer independent key management can connect external wallets to the same quote market, and embedded-wallet keys can be exported.

\subsection{Residual risks}\label{sec:risk}

Two exposures sit outside the enforcement boundary. Cross-chain settlement depends on filler liveness: market movement or unavailable fillers cause expiry and require a new quote, with multiple filler networks and alternative routes reducing this dependence where available. Deposit-address settlement depends on the venue delivering output or honoring a refund during the settlement window disclosed in the quote.

% !TEX root = ../main.tex
\section{Compliance Posture}\label{sec:compliance}\label{sec:credential}

Verification is optional. The permissionless lane requires none; quotes from regulated venues carry the provider's requirements, and eligibility, including jurisdictional restrictions, is checked per venue and per quote and disclosed before execution. A verified user holds a reusable network-level credential, reducing repeated onboarding across venues that accept it. Identity documents remain with the verification provider and the venue's own onboarding process; Pennysia coordinates only credential status.

Regulated coverage expands as venues are onboarded and may remain limited or unavailable in some jurisdictions. Pennysia enforces provider eligibility rules but does not determine a user's legal or tax status. Settlement through regulated venue rails carries the counterparty exposure described in Section~\ref{sec:risk}.

% !TEX root = ../main.tex
\section{Roadmap}\label{sec:roadmap}

Phases are ordered by dependency, not date. Each ships a working product (Table~\ref{tab:roadmap}).

\begin{table}[ht]
\centering\small
\caption{Rollout phases: ordered by dependency, each shipping a working product with a concrete exit criterion.}
\label{tab:roadmap}
\begin{tabular}{@{}>{\raggedright\arraybackslash}p{0.17\textwidth}p{0.44\textwidth}p{0.31\textwidth}@{}}
\toprule
\textbf{Phase} & \textbf{Ships} & \textbf{Exit when} \\
\midrule
\textbf{0. Prototype} \emph{(Ethereum mainnet)} &
Parallel venue competition, Settlement with atomic fallback, WAP floor, four audit toolchains, public docs/SDK/status surface &
The core loop (broadcast, compare, protect, execute) runs in production \\
\addlinespace[0.4em]
\textbf{1. Account + multi-chain spot} &
Embedded wallets (Web2 login $\rightarrow$ non-custodial EVM and Solana addresses), external-wallet parity, quote market across major EVM L2s, Solana, BNB Chain; signature-first same-chain fills &
A user with only a Web2 login and USDC executes a same-chain swap: no seed phrase, no native gas token \\
\addlinespace[0.4em]
\textbf{2. Chainless execution} &
Cross-chain parity: intent rails under the same one-signature flow, unified balance, Settlement on major EVM networks, fallback across route chains, deeper venue coverage &
Any asset on any chain to any asset on any chain, in one signature, WAP floor, one status line \\
\addlinespace[0.4em]
\textbf{3. Regulated lane} &
Verification tier; first centralized venues join the quote market; fiat ramps connect through the same credential (cards, bank, mobile pay); venue rails reach production; jurisdiction-aware eligibility &
A verified user's large order is won by a regulated venue, in an interface unverified users never notice missing \\
\addlinespace[0.4em]
\textbf{4. Network scale} &
New service verticals join the quote market (perpetuals, prediction markets, yield vaults); more chain families; public order-format spec; documented APIs/SDKs &
Third-party fillers and integrators build on the network \\
\bottomrule
\end{tabular}
\end{table}

This paper should be judged on four metrics: onboarding completion (account to first settled trade without a seed phrase or gas token), cross-chain fill quality versus the best single-chain alternative, venue count and share of fills won per lane, and fill-verification rate together with failure-mode frequency.


\FloatBarrier
\appendix
% !TEX root = ../main.tex
\section{The Ethereum Prototype}\label{app:v1}

This appendix records the experimental deployment on Ethereum Mainnet, live since August 28, 2026,\footnote{Live at \url{https://app.pennysia.com}; launch announcement at \url{https://x.com/PennysiaLabs/status/2093384908354048180}; contracts and verification reports at \url{https://github.com/Pennysia/settlement}.} on which the OLN's design is based.

\begin{itemize}[leftmargin=1.5em]
  \item \textbf{Parallel venue competition.} 30 venues quote per swap across four service classes (Table~\ref{tab:venues}), reaching thousands of underlying liquidity sources. 12,000+ tokens listed.
  \item \textbf{Settlement contract.} Publicly available onchain router under BUSL-1.1 (\texttt{0x3Aad\ldots795E}): single token pull, ordered venue fallback in one transaction, cross-chain intent routes, full revert on all-routes-fail.
  \item \textbf{Protections.} Signed output floor, gas-aware ranking, and unified Settlement approval, forming the basis of Section~\ref{sec:security}.
  \item \textbf{Fees live.} Surplus fee on outperformance (10\% cap, hardcoded), flat per-order protocol fee, partner share on intent routes. Every fee folds into the signed quote.
  \item \textbf{Automated verification.} Slither (static analysis), Halmos (symbolic execution, 168 tests), Medusa (coverage-guided fuzzing, 223 tests), Certora (formal verification, 105 rules). The latter three use isolated component harnesses. Reports are public in the Settlement repository.
  \item \textbf{Defensive invariants.} Reentrancy guards via EIP-1153~\cite{eip1153}; venue/allowance-target may not equal the contract. Refunds are baseline-relative. Owner keys cannot raise the surplus cap.
\end{itemize}

\begin{table}[ht]
\centering\small
\caption{The prototype's 30-venue coverage by service class.}
\label{tab:venues}
\begin{tabular}{@{}ll@{}}
\toprule
\textbf{Class} & \textbf{Venues} \\
\midrule
DEXs \& aggregators & onchain spot venues and meta-aggregators \\
Intent venues & same-chain and cross-chain intent networks \\
RFQ makers & quote-driven market makers \\
Privacy swaps & no-account, deposit-address providers \\
\bottomrule
\end{tabular}
\end{table}

% !TEX root = ../main.tex
% Keep the short schema and its explanation together.
\noindent\begin{minipage}{\textwidth}
\section{Order Format}\label{app:order}

The following payload describes the proposed order fields, using EIP-712 typed data on EVM~\cite{eip712}. The final specification will define encoding and chain-specific signatures.

{\footnotesize
\begin{verbatim}
OLN-Order {
  version:      uint16        // format version
  signer:       address       // embedded or external wallet
  input: {
    chain:      chainId       // origin chain
    token:      address       // asset sold
    amount:     uint256       // exact input
  }
  output: {
    chain:      chainId       // destination chain
    token:      address       // asset bought
    minAmount:  uint256       // WAP floor - worst acceptable output
  }
  lanes:        bitmask       // PERMISSIONLESS | REGULATED
  recipient:    address       // output delivery address
  deadline:     timestamp     // order expiry
  nonce:        uint256       // replay protection + uniqueness
  signature:    bytes         // EIP-712 typed-data (EVM) / equiv. (Solana)
}
\end{verbatim}
}

The output floor \texttt{minAmount} is the WAP. The \texttt{lanes} field limits eligible venues, with \texttt{REGULATED} requiring verification. Each rail uses these fields to validate delivery to the specified recipient on the destination chain before the deadline.
\end{minipage}\par


% !TEX root = ../main.tex
{\footnotesize
\begin{thebibliography}{99}
\setlength{\itemsep}{0pt}\setlength{\parskip}{0pt}

\bibitem{eip1153}
A.~Akhunov and M.~Salem,
``EIP-1153: Transient storage opcodes,'' Ethereum Improvement Proposal, 2018.
\url{https://eips.ethereum.org/EIPS/eip-1153}

\bibitem{eip712}
R.~Bloemen, L.~Logvinov, and J.~Evans,
``EIP-712: Typed structured data hashing and signing,'' Ethereum Improvement Proposal, 2017.
\url{https://eips.ethereum.org/EIPS/eip-712}

\bibitem{erc4337}
V.~Buterin et al.,
``ERC-4337: Account abstraction using alt mempool,'' Ethereum Improvement Proposal, 2021.
\url{https://eips.ethereum.org/EIPS/eip-4337}

\bibitem{eip7702}
V.~Buterin et al.,
``EIP-7702: Set code for EOAs,'' Ethereum Improvement Proposal, 2024.
\url{https://eips.ethereum.org/EIPS/eip-7702}

\bibitem{intentauction}
A.~Canidio and F.~Henneke,
``Fair Combinatorial Auction for Blockchain Trade Intents: Being Fair without Knowing What is Fair,'' 2025. arXiv:2408.12225.
\url{https://arxiv.org/abs/2408.12225}

\bibitem{smartaccounts}
Castle Labs,
``Smart Accounts: A New Era for UX,'' research report, April 2025.
\url{https://castlelabs.io/research/smart-accounts-a-new-era-for-ux}

\bibitem{coingecko}
CoinGecko Research,
``DEX to CEX volume ratios reach new highs in 2025,'' research report, CoinGecko, 2025.
\url{https://www.coingecko.com/research/publications/dex-to-cex-ratio}

\bibitem{flashboys}
P.~Daian, S.~Goldfeder, T.~Kell, Y.~Li, X.~Zhao, I.~Bentov, L.~Breidenbach, and A.~Juels,
``Flash Boys 2.0: Frontrunning, Transaction Reordering, and Consensus Instability in Decentralized Exchanges,'' \emph{Proc.\ IEEE S\&P}, 2020. arXiv:1904.05234.
\url{https://arxiv.org/abs/1904.05234}

\bibitem{threshold}
R.~Gennaro and S.~Goldfeder,
``Fast Multiparty Threshold ECDSA with Fast Trustless Setup,'' \emph{Proc.\ ACM CCS}, 2018.
\url{https://eprint.iacr.org/2019/114}

\bibitem{atomicswaps}
M.~Herlihy,
``Atomic Cross-Chain Swaps,'' \emph{Proc.\ ACM PODC}, 2018. arXiv:1801.09515.
\url{https://arxiv.org/abs/1801.09515}

\bibitem{bitcoin}
S.~Nakamoto,
``Bitcoin: A Peer-to-Peer Electronic Cash System,'' 2008.
\url{https://bitcoin.org/bitcoin.pdf}

\bibitem{theblock}
B.~Samsoedin,
``Data \& Insights: DEX to CEX ratio Skyrockets; Polymarket US Volumes Too,'' \emph{The Block}, October 7, 2026.
\url{https://www.theblock.co/newsletters/data-and-insights/2026-10-07-data-polymarket-us-volumes-too-417735}

\bibitem{sok-defi}
S.~M. Werner, D.~Perez, L.~Gudgeon, A.~Klages-Mundt, D.~Harz, and W.~J. Knottenbelt,
``SoK: Decentralized Finance (DeFi),'' \emph{Proc.\ ACM AFT}, 2022. arXiv:2101.08778.
\url{https://arxiv.org/abs/2101.08778}

\bibitem{solverwelfare}
Y.~Yuminaga, D.~Chen, and D.~Sui,
``Execution Welfare Across Solver-based DEXes,'' 2025. arXiv:2503.00738.
\url{https://arxiv.org/abs/2503.00738}

\bibitem{bridges}
M.~Zhang, X.~Zhang, J.~Barbee, Y.~Zhang, and Z.~Lin,
``SoK: Security of Cross-chain Bridges: Attack Surfaces, Defenses, and Open Problems,'' arXiv preprint arXiv:2312.12573, 2023.
\url{https://arxiv.org/abs/2312.12573}

\bibitem{dexdepth}
B.~Zhu, D.~Liu, X.~Wan, G.~Liao, C.~Moallemi, and B.~Bachu,
``What Drives Liquidity on Decentralized Exchanges? Evidence from the Uniswap Protocol,'' \emph{Proc.\ CAAW}, 2025. arXiv:2410.19107.
\url{https://arxiv.org/abs/2410.19107}

\end{thebibliography}
}


\end{document}
